\chapter{Conclusion et perspectives}
\label{ch:conclusion}

\section{Bilan}

Ce stage à WEBVOICE a permis de livrer \textbf{LexHub}, plateforme d'assistance IA pour cabinets d'avocats, conformément au sujet du PFA. L'objectif n'était pas un démonstrateur hors-ligne : c'est un produit déployé sur \url{https://lexhub.space}, avec isolation multi-tenant, rôles, atelier de plaidoirie guidé, validation / assignation, et export officiel.

Sur le plan académique, le travail relie analyse des besoins, étude comparative, conception UML, réalisation full-stack et mise en production. Sur le plan professionnel, il confronte aux contraintes réelles : e-mails, temps réel, export, déontologie, robustesse des dépendances IA.

\section{Apports personnels}

Maîtrise d'une architecture web moderne, conception d'un flux métier (voix $\rightarrow$ acte), vigilance vis-à-vis des inventions de l'IA, dialogue avec un tuteur école (M.~Mohammed Al Achhab) et un encadrant entreprise (M.~Hatim Benhaddou).

\section{Perspectives}

\begin{itemize}[leftmargin=1.4em]
  \item Amélioration continue de la transcription (accents, conditions acoustiques) et des alternatives STT.
  \item Consolidation du lien planning $\leftrightarrow$ dossiers et des filtres métier.
  \item Signature électronique et parcours de validation multi-niveaux.
  \item Facturation SaaS et tableaux de bord admin enrichis.
  \item CI/CD et durcissement sécurité avancé --- sans sortir du cœur plaidoirie.
\end{itemize}
